\documentclass[12pt, a4paper]{ctexart}
\usepackage{fancyhdr}
\usepackage{amsmath, amssymb}
\usepackage{geometry}
\usepackage{enumitem}
\usepackage{setspace}
\usepackage{graphicx}
\usepackage{tabularx}

\geometry{left=2cm, right=2cm, top=2.5cm, bottom=2.5cm}
\setCJKmainfont{Hiragino Sans GB}
\setCJKsansfont{Heiti SC}
\setstretch{1.5}

\pagestyle{fancy}
\fancyhf{}
\fancyhead[L]{}
\fancyhead[R]{\small 数学（选择性必修第三册）}
\fancyfoot[C]{\thepage}
\renewcommand{\headrulewidth}{0.4pt}

\newcommand{\score}[1]{\hfill\textbf{（#1分）}}

\begin{document}

\begin{center}
\Large\textbf{选择性必修第三册·第六章 计数原理 测试卷}

\vspace{0.3cm}
\normalsize 满分：150分 \quad 考试时间：120分钟
\end{center}

\vspace{0.5cm}

\noindent\textbf{一、单项选择题（本大题共8小题，每小题5分，共40分。每小题只有一个选项符合题目要求。）}

\vspace{0.3cm}

\begin{enumerate}[label=\textbf{\arabic*.}, leftmargin=2em]

\item 从5名同学中选3人参加社区志愿服务活动，不同的选法有\score{5}
\begin{enumerate}[label=\Alph*., leftmargin=2em]
\item 10种 \quad B. 15种 \quad C. 20种 \quad D. 60种
\end{enumerate}

\item 用0，1，2，3，4，5这6个数字，可以组成没有重复数字的三位偶数的个数是\score{5}
\begin{enumerate}[label=\Alph*., leftmargin=2em]
\item 52 \quad B. 60 \quad C. 68 \quad D. 72
\end{enumerate}

\item $\left(x-\dfrac{1}{x}\right)^6$的展开式中含$x^2$项的系数是\score{5}
\begin{enumerate}[label=\Alph*., leftmargin=2em]
\item $-20$ \quad B. $-15$ \quad C. $15$ \quad D. $20$
\end{enumerate}

\item 将4本不同的书全部分给甲、乙、丙三人，每人至少1本，则不同的分法有\score{5}
\begin{enumerate}[label=\Alph*., leftmargin=2em]
\item 12种 \quad B. 24种 \quad C. 36种 \quad D. 48种
\end{enumerate}

\item 同时抛掷两枚骰子，向上的点数之和为7的概率是\score{5}
\begin{enumerate}[label=\Alph*., leftmargin=2em]
\item $\dfrac{1}{6}$ \quad B. $\dfrac{1}{8}$ \quad C. $\dfrac{5}{36}$ \quad D. $\dfrac{7}{36}$
\end{enumerate}

\item 二项式$(x+1)^8$的展开式中$x^5$的系数是\score{5}
\begin{enumerate}[label=\Alph*., leftmargin=2em]
\item 28 \quad B. 56 \quad C. 70 \quad D. 112
\end{enumerate}

\item 6名同学站成一排照相，其中甲、乙两人必须相邻的排法有\score{5}
\begin{enumerate}[label=\Alph*., leftmargin=2em]
\item 120种 \quad B. 240种 \quad C. 360种 \quad D. 480种
\end{enumerate}

\item 将5封不同的信投入3个不同的信箱，每个信箱至少投1封信，则不同的投法有\score{5}
\begin{enumerate}[label=\Alph*., leftmargin=2em]
\item 120种 \quad B. 150种 \quad C. 180种 \quad D. 240种
\end{enumerate}

\end{enumerate}

\vspace{0.5cm}

\noindent\textbf{二、多项选择题（本大题共4小题，每小题5分，共20分。每小题有多个选项符合题目要求，全部选对得5分，部分选对得2分，有选错的得0分。）}

\vspace{0.3cm}

\begin{enumerate}[label=\textbf{\arabic*.}, leftmargin=2em, resume]

\item 下列等式正确的是\score{5}
\begin{enumerate}[label=\Alph*., leftmargin=2em]
\item $C_n^m = C_n^{n-m}$
\item $A_n^m = n \cdot A_{n-1}^{m-1}$
\item $C_n^m = C_{n-1}^m + C_{n-1}^{m-1}$
\item $A_n^m = C_n^m \cdot m!$
\end{enumerate}

\item 关于二项式$(1+x)^{10}$，下列说法正确的是\score{5}
\begin{enumerate}[label=\Alph*., leftmargin=2em]
\item 展开式共有11项
\item 展开式中各项系数之和为$2^{10}$
\item 展开式中$x^4$的系数与$x^6$的系数相等
\item 展开式中二项式系数最大的项是第6项
\end{enumerate}

\item 从5双不同的鞋子中任取4只，则下列说法正确的是\score{5}
\begin{enumerate}[label=\Alph*., leftmargin=2em]
\item 恰好有2只配成一双的取法有120种
\item 恰好有2只配成一双的取法有140种
\item 4只都不成双的取法有80种
\item 至少有2只成双的取法有130种
\end{enumerate}

\item 将4个不同的小球放入3个不同的盒子中，下列说法正确的是\score{5}
\begin{enumerate}[label=\Alph*., leftmargin=2em]
\item 不同的放法共有$3^4$种
\item 恰有一个空盒的放法有42种
\item 每个盒子都不空的放法有36种
\item 恰好有两个盒子为空的放法有6种
\end{enumerate}

\end{enumerate}

\vspace{0.5cm}

\noindent\textbf{三、填空题（本大题共4小题，每小题5分，共20分。）}

\vspace{0.3cm}

\begin{enumerate}[label=\textbf{\arabic*.}, leftmargin=2em, resume]

\item 二项式$\left(x^2+\dfrac{1}{x}\right)^6$的展开式中的常数项为\hspace{3cm}。\score{5}

\item 已知$C_{10}^x = C_{10}^3$，则$x=$\hspace{3cm}。\score{5}

\item 二项式$(2x-1)^5$的展开式中各项系数之和为\hspace{3cm}。\score{5}

\item 从1，2，3，4，5这5个数字中任取3个不同的数字排成三位数，其中百位数字比十位数字大，十位数字比个位数字大的三位数共有\hspace{3cm}个。\score{5}

\end{enumerate}

\newpage

\noindent\textbf{四、解答题（本大题共6小题，共70分。解答应写出必要的文字说明、证明过程或演算步骤。）}

\vspace{0.3cm}

\begin{enumerate}[label=\textbf{\arabic*.}, leftmargin=2em, resume]

\item（10分）某班有4名男生和3名女生。\score{10}

（1）从中选3人参加数学竞赛，要求至少有1名女生，有多少种选法？

（2）从中选3人分别参加数学、物理、化学竞赛，要求至少有1名女生，有多少种选法？

\item（12分）已知$(2x-1)^6 = a_0 + a_1 x + a_2 x^2 + a_3 x^3 + a_4 x^4 + a_5 x^5 + a_6 x^6$。\score{12}

（1）求$a_0$的值；

（2）求$a_1 + a_2 + a_3 + a_4 + a_5 + a_6$的值；

（3）求$a_0 + a_2 + a_4 + a_6$的值。

\item（12分）一个口袋中装有4个红球和3个白球，这些球除颜色外完全相同。\score{12}

（1）从袋中一次取出3个球，求恰好取到2个红球的概率；

（2）从袋中依次不放回地取3个球，求第3次取到白球的概率。

\item（12分）用0，1，2，3，4，5这6个数字组成没有重复数字的数。\score{12}

（1）能组成多少个没有重复数字的四位偶数？

（2）能组成多少个没有重复数字且比3400大的四位数？

\item（12分）已知$(x^2+ax)^5$的展开式中$x^6$的系数为80。\score{12}

（1）求实数$a$的值；

（2）求展开式中$x^4$的系数；

（3）求展开式中二项式系数最大的项。

\item（12分）证明题。\score{12}

（1）用二项式定理证明：$2^n > n^2$（其中$n \geqslant 5$，$n \in \mathbb{N}^*$）；

（2）用二项式定理证明：$3^{2n+2} - 8n - 9$能被64整除（其中$n \in \mathbb{N}^*$）。

\end{enumerate}

\vspace{1cm}
\begin{center}
\rule{8cm}{0.4pt}

\vspace{0.3cm}
\textbf{本试卷结束}
\end{center}

\end{document}
